\chapter{การจัดการระบบน้ำอัจฉริยะและการควบคุมการชลประทานแบบแม่นยำสูง}
\label{chap:ch04}

\section*{แผนบริหารการสอนประจำบทที่ 4}
\addcontentsline{toc}{section}{แผนบริหารการสอนประจำบทที่ 4}

\noindent\textbf{หัวข้อเนื้อหาประจำบท}
\begin{enumerate}
    \item สรีรวิทยาความต้องการน้ำของทุเรียนในแต่ละระยะฟีโนโลยี
    \item การคำนวณการใช้น้ำของพืชตามแบบจำลอง
    \item สถาปัตยกรรมระบบควบคุมน้ำอัตโนมัติด้วยโซลินอยด์วาล์ว
    \item การประหยัดน้ำและพลังงานด้วยเทคโนโลยีอัตโนมัติ
\end{enumerate}

\noindent\textbf{วัตถุประสงค์เชิงพฤติกรรม}
\begin{enumerate}
    \item อธิบายหลักการ ทฤษฎี และสถาปัตยกรรมเทคโนโลยีประจำบทที่ 4 ได้อย่างถูกต้อง
    \item ประยุกต์ใช้เครื่องมือ อุปกรณ์เซนเซอร์ และระบบปัญญาประดิษฐ์เพื่อการจัดการสวนทุเรียนได้อย่างมีประสิทธิภาพ
    \item วิเคราะห์ สรุปผลข้อมูล และแก้ไขปัญหาเชิงฟิสิกส์เกษตรและคุณภาพดินได้อย่างเป็นระบบ
\end{enumerate}

\noindent\textbf{การวัดและประเมินผล}
\begin{table}[H]
\centering\small
\begin{tabularx}{\textwidth}{p{0.55\textwidth} p{0.25\textwidth} c}
\toprule
\textbf{เกณฑ์การประเมินผล} & \textbf{เครื่องมือวัดผล} & \textbf{สัดส่วนคะแนน} \\
\midrule
1. ทักษะการปฏิบัติงานและการใช้อุปกรณ์ & แบบประเมินทักษะภาคสนาม & 40\% \\
2. รายงานข้อมูลและการวิเคราะห์ผล & การตรวจสมุดบันทึกและดาต้าชีท & 30\% \\
3. แบบฝึกหัดและคำถามท้ายบท & แบบทดสอบท้ายบทเรียน & 20\% \\
4. การมีส่วนร่วมและความรับผิดชอบ & การเข้าเรียนและการตรงต่อเวลา & 10\% \\
\midrule
\textbf{รวมทั้งสิ้น} & & \textbf{100\%} \\
\bottomrule
\end{tabularx}
\end{table}
\newpage

% ==========================================================
% เนื้อหาบทเรียนประจำบทที่ 4
% ==========================================================

\section{สรีรวิทยาความต้องการน้ำของทุเรียนในแต่ละระยะฟีโนโลยี}
\index{สรีรวิทยาความต้องการน้ำของทุเรียนในแต่ละระยะฟีโนโลยี}

การศึกษาและทำความเข้าใจเกี่ยวกับสรีรวิทยาความต้องการน้ำของทุเรียนในแต่ละระยะฟีโนโลยี (Durian Water Requirements Across Phenological Stages) มีความสำคัญอย่างยิ่งในระบบเกษตรอัจฉริยะ โดยมีรายละเอียดสาระสำคัญดังนี้

ทุเรียนเป็นไม้ผลที่มีความต้องการน้ำในปริมาณและจังหวะเวลาที่จำเพาะเจาะจงอย่างยิ่งในแต่ละระยะการเจริญเติบโต (Phenological stages):

1. ระยะสะสมอาหารและเตรียมตาดอก (Resting \& Floral induction): ต้องงดการให้น้ำ (กักน้ำ) เป็นเวลา 10-14 วัน จนกระทั่งดินแห้งสนิท ใบมีอาการสลดเล็กน้อย เพื่อกระตุ้นให้ต้นทุเรียนเกิดความเครียดจากการขาดน้ำ (Water stress) ผลักดันให้พัฒนาตาดอก
2. ระยะพัฒนาตาดอกถึงเหยียดตีนเป็ด (Flower development): เริ่มให้น้ำปริมาณน้อยแบบโชยๆ ประมาณ 1 ใน 3 ของปกติ แล้วค่อยๆ เพิ่มปริมาณน้ำอย่างช้าๆ
3. ระยะดอกบาน (Anthesis): ลดปริมาณน้ำลงเหลือประมาณ 50\% เพื่อป้องกันการสลัดดอกและช่วยให้เกสรผสมติดสมบูรณ์
4. ระยะพัฒนาการของผล (Fruit expansion): ทุเรียนต้องการน้ำในปริมาณสูงและสม่ำเสมอที่สุด การขาดน้ำในระยะนี้จะทำให้ผลแคระแกร็น หนามบิดเบี้ยว และสลัดผลทิ้ง
5. ระยะก่อนเก็บเกี่ยว 10-14 วัน: ลดปริมาณน้ำลงเล็กน้อยเพื่อช่วยให้เนื้อทุเรียนสะสมแป้งและน้ำตาลเต็มที่ ป้องกันอาการเนื้อแฉะเต่าเผา

\section{การคำนวณการใช้น้ำของพืชตามแบบจำลอง}
\index{การคำนวณการใช้น้ำของพืชตามแบบจำลอง}

การศึกษาและทำความเข้าใจเกี่ยวกับการคำนวณการใช้น้ำของพืชตามแบบจำลอง (Crop Evapotranspiration and Water Budgeting) มีความสำคัญอย่างยิ่งในระบบเกษตรอัจฉริยะ โดยมีรายละเอียดสาระสำคัญดังนี้

การคำนวณความต้องการน้ำของต้นทุเรียนต่อวันอาศัยสมการมาตรฐานของ FAO-56:
\[ ET\_c = K\_c \times ET\_0 \]
โดยที่ $ET_c$ คือการคายระเหยน้ำของพืช (Crop evapotranspiration, mm/day), $K_c$ คือสัมประสิทธิ์พืชทุเรียน (Crop coefficient ซึ่งแปรผันตามระยะการเจริญเติบโต ตั้งแต่ 0.4 ถึง 0.85), และ $ET_0$ คือการคายระเหยน้ำอ้างอิงจากสถานีตรวจอากาศ คำนวณด้วยสมการ Penman-Monteith

ปริมาตรน้ำที่ต้องให้ต่อต้นต่อวัน ($V$, ลิตร/ต้น/วัน) คำนวณได้จาก:
\[ V = ET\_c \times A\_{\text{canopy}} \times \frac{1}{E\_i} \]
โดยที่ $A_{\text{canopy}}$ คือพื้นที่ทรงพุ่มทุเรียน ($\pi r^2$, ตารางเมตร) และ $E_i$ คือประสิทธิภาพของระบบการให้น้ำแบบมินิสปริงเกลอร์ (ประมาณ 85-90\%)

\section{สถาปัตยกรรมระบบควบคุมน้ำอัตโนมัติด้วยโซลินอยด์วาล์ว}
\index{สถาปัตยกรรมระบบควบคุมน้ำอัตโนมัติด้วยโซลินอยด์วาล์ว}

การศึกษาและทำความเข้าใจเกี่ยวกับสถาปัตยกรรมระบบควบคุมน้ำอัตโนมัติด้วยโซลินอยด์วาล์ว (Automated Solenoid Valve Control Architecture) มีความสำคัญอย่างยิ่งในระบบเกษตรอัจฉริยะ โดยมีรายละเอียดสาระสำคัญดังนี้

ระบบชลประทานอัจฉริยะผสานไมโครคอนโทรลเลอร์ควบคุมหลัก (เช่น ESP32 หรือ Industrial PLC) ร่วมกับบอร์ดขับรีเลย์เพื่อสั่งการเปิด-ปิด โซลินอยด์วาล์วไฟฟ้า (Solenoid Valve 24VAC หรือ 12VDC Latching)

โหมดการทำงานของระบบ:
1. การควบคุมตามเวลา (Time-based scheduling): สั่งเปิด-ปิดน้ำตามกำหนดเวลาที่ตั้งล่วงหน้า
2. การควบคุมตามค่าเซนเซอร์ความชื้นในดิน (Sensor-based thresholding): สั่งเปิดน้ำอัตโนมัติเมื่อความชื้นในดินลดต่ำลงถึงจุดวิกฤต (Lower threshold) และหยุดให้น้ำเมื่อดินมีความชื้นถึงระดับความจุภาคสนาม (Field Capacity)
3. การควบคุมอิงการคำนวณสมดุลน้ำ (Evapotranspiration-based precision): ประมวลผลร่วมกับสถานีตรวจอากาศเพื่อให้น้ำชดเชยตามปริมาณน้ำที่สูญเสียไปจริงในแต่ละวัน

\begin{techbox}[ข้อกำหนดทางเทคนิคสถาปัตยกรรม AIoT]
การเชื่อมต่อระบบเซนเซอร์ไร้สายต้องรักษาความมั่นคงปลอดภัยของข้อมูล มีการเข้ารหัสระดับ AES-128 และบริหารจัดการพลังงานให้โหนดเซนเซอร์ทำงานได้อย่างต่อเนื่องไม่น้อยกว่า 3 ปี
\end{techbox}

\section{การประหยัดน้ำและพลังงานด้วยเทคโนโลยีอัตโนมัติ}
\index{การประหยัดน้ำและพลังงานด้วยเทคโนโลยีอัตโนมัติ}

การศึกษาและทำความเข้าใจเกี่ยวกับการประหยัดน้ำและพลังงานด้วยเทคโนโลยีอัตโนมัติ (Water and Energy Efficiency Optimization) มีความสำคัญอย่างยิ่งในระบบเกษตรอัจฉริยะ โดยมีรายละเอียดสาระสำคัญดังนี้

จากการทดสอบภาคสนามในสวนทุเรียนจังหวัดจันทบุรี การเปลี่ยนจากการให้น้ำตามความเคยชินของแรงงานมาเป็นระบบอัตโนมัติแบบแม่นยำสูง สามารถลดปริมาณการใช้น้ำลงได้ 25 ถึง 35\% พร้อมทั้งลดชั่วโมงการทำงานของเครื่องสูบน้ำ ส่งผลให้ประหยัดพลังงานไฟฟ้าและน้ำมันเชื้อเพลิงได้มากกว่า 30\% ในขณะที่คุณภาพผลผลิตทุเรียนสม่ำเสมอเพิ่มขึ้น

\section{ขั้นตอนและวิธีปฏิบัติการทดลอง}
\index{ขั้นตอนการทดลองประจำบทที่ 4}

\subsection{การคำนวณความต้องการน้ำของต้นทุเรียนอายุ 8 ปี}
\begin{sensorbox}[การคำนวณความต้องการน้ำของต้นทุเรียนอายุ 8 ปี]
1. วัดรัศมีทรงพุ่มทุเรียนเฉลี่ย $r = 3.5$ เมตร คำนวณพื้นที่ทรงพุ่ม: $A = \pi \times (3.5)^2 \approx 38.5$ ตารางเมตร
2. อ่านค่า $ET_0$ จากสถานีอากาศ AIoT ในวันแดดจัด ได้ค่า $ET_0 = 5.2$ mm/day
3. ในระยะผลขยายขนาด ค่าสัมประสิทธิ์พืช $K_c = 0.80$ คำนวณ: $ET_c = 0.80 \times 5.2 = 4.16$ mm/day
4. คำนวณปริมาตรน้ำที่ต้องให้: $V = 4.16 \times 38.5 / 0.85 \approx 188.4$ ลิตร/ต้น/วัน
5. กำหนดเวลาเปิดสปริงเกลอร์: หากมินิสปริงเกลอร์มีอัตราจ่ายน้ำรวม 120 ลิตร/ชั่วโมง ต้องเปิดน้ำเป็นเวลา 1 ชั่วโมง 34 นาที
\end{sensorbox}

\subsection{การติดตั้งและทดสอบโซลินอยด์วาล์วไฟฟ้า 24VAC ประจำแปลง}
\begin{sensorbox}[การติดตั้งและทดสอบโซลินอยด์วาล์วไฟฟ้า 24VAC ประจำแปลง]
1. ติดตั้งโซลินอยด์วาล์วขนาด 2 นิ้วที่ท่อแยกเข้าสู่แต่ละโซนรดน้ำ หุ้มกล่องวาล์วกันดิน (Valve Box)
2. เดินสายไฟควบคุมทนน้ำ (UF Cable) ฝังดินร่วมกับท่อประปาเข้าสู่ตู้คอนโทรลเลอร์
3. ทดสอบแรงดันน้ำก่อนเข้าและหลังผ่านวาล์ว ด้วยเกจวัดแรงดัน (Pressure Gauge) ให้อยู่ในช่วง 1.5 - 2.5 บาร์
4. ทดสอบสั่งเปิด-ปิดวาล์วผ่านแอปพลิเคชันสมาร์ทโฟนและตรวจสอบกระแสไฟกระตุ้นคอยล์วาล์ว
\end{sensorbox}

\subsection{การเขียนฟังก์ชันป้องกันความผิดพลาดของระบบน้ำ}
\begin{sensorbox}[การเขียนฟังก์ชันป้องกันความผิดพลาดของระบบน้ำ]
1. ติดตั้งเซนเซอร์วัดการไหลของน้ำ (Water Flow Meter) ที่ท่อเมนหลัก
2. กำหนดเงื่อนไข Failsafe: หากระบบสั่งเปิดวาล์วแต่ Flow Meter ตรวจไม่พบการไหลภายใน 60 วินาที ให้สั่งตัดปั๊มน้ำทันทีเพื่อป้องกันปั๊มไหม้ (Dry-run protection)
3. หากตรวจพบอัตราการไหลสูงเกินปกติ 150\% ให้แจ้งเตือนท่อเมนแตก (Pipe burst alert)
\end{sensorbox}

\section{ตารางบันทึกผลการทดลองและการอภิปรายผล}
\begin{table}[H]
\centering\small
\caption{ตารางบันทึกผลการตรวจวัดและข้อมูลพารามิเตอร์ประจำบทที่ 4}
\label{tab:ch04_datasheet}
\begin{tabularx}{\textwidth}{p{0.3\textwidth} p{0.3\textwidth} X}
\toprule
\textbf{พารามิเตอร์ / การทดสอบ} & \textbf{ค่าที่ตรวจวัดได้} & \textbf{การแปลผลและการวิเคราะห์} \\
\midrule
จุดตรวจวัดที่ 1 (บริเวณดินดอน) & ค่าความชื้นอยู่ในเกณฑ์ปกติ & ปริมาณน้ำเพียงพอต่อการเจริญ \\
จุดตรวจวัดที่ 2 (บริเวณที่ลุ่ม) & มีการสะสมความชื้นสูง & ควรชะลอการให้น้ำเพื่อป้องกันรากเน่า \\
สถานะสัญญาณโครงข่าย & RSSI: -85 dBm, SNR: +8 dB & คุณภาพสัญญาณ LoRaWAN ยอดเยี่ยม \\
\bottomrule
\end{tabularx}
\end{table}

\section{คำถามทบทวนและแบบฝึกหัดท้ายบท}
\begin{enumerate}
    \item จงอธิบายเหตุผลทางสรีรวิทยาพืชว่าเหตุใดจึงต้องงดการให้น้ำทุเรียนเพื่อชักนำการออกดอก
    \item หากเกิดฝนตกหนักวัดปริมาณได้ 35 mm ในช่วงเช้า อัลกอริทึมชลประทานอัจฉริยะควรปรับแผนการให้น้ำในวันนั้นอย่างไร
    \item จงเปรียบเทียบการให้น้ำแบบครั้งเดียวต่อวันกับการแบ่งให้น้ำเป็น 2 ช่วง (เช้าและบ่าย) ในแง่การดูดซึมของรากและการระเหยสูญเสีย
\end{enumerate}
